% ============================================================
%  When Synthetic Data Falls Short
%  ACL 2026 Short Paper Submission
%  Build:  pdflatex main.tex  (twice for cross-references)
% ============================================================
\documentclass[11pt]{article}
\usepackage{natbib} 
\usepackage{times}
\usepackage{latexsym}
\usepackage{microtype}
\usepackage{graphicx}
\usepackage{booktabs}
\usepackage{amsmath}
\usepackage{multirow}
\usepackage{xcolor}
\usepackage{url}

\renewcommand{\UrlFont}{\ttfamily\small}

% ── Metadata ─────────────────────────────────────────────────
\title{When Synthetic Data Falls Short:\\
A Pareto Analysis of LLM-Generated Data Filtering\\
for Hate Speech Detection}

\author{
  Anonymous Author(s)\\
  \textit{Anonymized for Review}
}

\begin{document}
\maketitle

% ── Abstract ─────────────────────────────────────────────────
\begin{abstract}
Annotation bottlenecks motivate the use of large language model (LLM) generated
synthetic data for low-resource hate speech detection.
We conduct a systematic Pareto analysis of discriminator-guided quality filtering
across six continuous thresholds using an open-source LLM (LLaMA 3.1-8B via Groq API),
the first such sweep for hate speech data.
Across all thresholds, synthetic-only training consistently underperforms real data:
the best synthetic condition achieves F1\,=\,0.576 versus 0.715 for real-only training,
a gap quality filtering cannot close.
A novel precision-recall asymmetry analysis reveals the diagnostic mechanism:
LLM-generated hate speech captures explicit, prototypical patterns well
(precision $\uparrow$) but systematically misses borderline and implicit cases
present in real annotation (recall $\downarrow$).
Augmentation with filtered synthetic data yields an F1 gain of $+0.012$
over real-only, which falls short of significance under McNemar's test
($\chi^2 = 3.125$, $p = 0.077$, $n = 400$).
Our findings urge practitioners to treat LLM-generated hate speech data
as supplementary rather than substitutional, and we release the full pipeline
for reproducibility.
\end{abstract}

% ── 1. Introduction ──────────────────────────────────────────
\section{Introduction}
\label{sec:intro}

Hate speech detection remains a high-stakes NLP task where annotated data is scarce,
expensive, and subject to annotation fatigue \cite{fortuna-nunes-2018-survey}.
The proliferation of capable open-source large language models (LLMs) has made
synthetic data generation an attractive alternative to human annotation,
with several studies reporting encouraging results for text classification
augmentation \cite{wang-etal-2021-want}.

The closest prior work to ours is \cite{casula-etal-2024-delving}, who apply
\emph{binary} quality filtering to GPT-generated synthetic data for Italian hate speech
and report mixed, task-dependent results.
However, three important gaps remain:
(1) no systematic \emph{continuous} threshold sweep has been conducted,
making it impossible to characterize the full quality-quantity Pareto frontier;
(2) GPT-based generation limits reproducibility; and
(3) the \emph{mechanism} behind synthetic data failures is unexplored.

We address all three gaps. Our contributions are:
\begin{enumerate}
  \item \textbf{First Pareto analysis} of the quality-quantity tradeoff for hate speech
    synthetic data, sweeping six quality thresholds continuously.
  \item \textbf{Precision-recall asymmetry} as a diagnostic tool: we show that
    LLM-generated data systematically biases classifiers toward explicit hate patterns,
    lowering recall on borderline cases.
  \item \textbf{Reproducible open-source pipeline} using LLaMA 3.1-8B via the Groq API
    and a discriminator-based quality proxy, enabling replication on any target dataset.
\end{enumerate}

% ── 2. Related Work ──────────────────────────────────────────
\section{Related Work}
\label{sec:related}

\paragraph{Synthetic data for NLP.}
Data augmentation via LLMs has shown benefits for low-resource classification, with Wang et al. [7] demonstrating that LLM-based annotation can reduce labeling costs, though domain-specific limits apply.
Our work provides the first systematic Pareto characterization for hate speech,
a domain with particularly high annotation sensitivity.

\paragraph{Hate speech detection and synthetic data.}
\cite{casula-etal-2024-delving} is the most direct predecessor: they generate synthetic Italian
hate speech with GPT models and apply binary quality filtering, finding inconsistent benefits.
\cite{fortuna-nunes-2018-survey} survey annotation challenges in multilingual settings,
while \cite{waseem-hovy-2016-hateful} established key benchmark datasets for English hate speech.
We extend this line by providing a continuous threshold analysis,
an open-source generator, and a diagnostic explanation for failure modes.

\paragraph{Data quality filtering.}
Quality-based selection is common in pre-training data curation
\cite{penedo-etal-2023-refinedweb,xie-etal-2023-data}.
\cite{xie-etal-2023-data} propose importance resampling to align synthetic distributions with target domains.
Unlike these works, which apply filtering to pre-training corpora, we apply it to
task-specific fine-tuning data and measure downstream classification performance
at a fine-grained threshold level.

% ── 3. Method ────────────────────────────────────────────────
\section{Method}
\label{sec:method}

\subsection{Synthetic Data Generation}
We generate 5{,}000 synthetic samples (2{,}500 hate, 2{,}500 non-hate) using
LLaMA 3.1-8B \cite{grattafiori-etal-2024-llama3} served via the Groq API.
Five prompt templates per class capture diverse stylistic registers:
professional misogyny, dehumanization, counter-speech, neutral discourse,
and empowerment framing.
Each seed is augmented with synonym swapping and prefix/suffix injection,
with a uniqueness check to suppress near-duplicates.
Samples with pairwise cosine similarity $\geq 0.95$ (sentence-BERT embeddings)
are removed; the surviving deduplicated set contains 1{,}170 samples.

\subsection{Discriminator Quality Proxy}
We train a binary classifier to distinguish real from synthetic text,
using sentence-BERT embeddings \cite{reimers-gurevych-2019-sentence} as features.
The \emph{quality score} is defined as:
\[
  q(x) = P(\text{real} \mid x)
\]
Samples with $q(x)$ near 0.5 are maximally ambiguous — indistinguishable
from real human-written text — and are thus considered highest quality.
A Random Forest discriminator (AUC\,=\,1.00 on a held-out set) provides the scores,
reflecting the substantial distributional gap between real and synthetic text.

\subsection{Continuous Threshold Filtering}
We define six conditions by retaining synthetic samples above quality percentile thresholds:
top-10\%, 30\%, 50\%, 70\%, 90\%, and unfiltered (100\%).
We additionally construct a \textsc{real+top50} augmented set combining all 500 real samples
with the top-50\% quality synthetic samples.
This yields eight distinct training conditions for the Pareto analysis.

% ── 4. Experiments ───────────────────────────────────────────
\section{Experiments}
\label{sec:experiments}

\subsection{Dataset and Setup}
We use the \textbf{HatEval} dataset \cite{basile-etal-2019-semeval},
a benchmark for misogyny and hate speech targeting women and immigrants in English.
We access the dataset via the \textbf{Lots-of-LoRAs/task333\_hateeval\_classification\_hate\_en}
split on HuggingFace, which provides a preprocessed English binary classification
version of the original SemEval 2019 Task 5 data.
We sample 500 training examples and 400 test examples (200 hate, 200 non-hate),
with zero overlap enforced by excluding test indices before sampling the training pool.
The 400-example test set provides sufficient statistical power to detect medium effects
($d \geq 0.3$) at conventional levels ($\alpha = 0.05$, $\beta = 0.80$).
All conditions are evaluated on the same held-out test set.

The classifier is \textbf{MiniLM-L12-H384} \cite{wang-etal-2020-minilm}, fine-tuned for 3 epochs with $\mathrm{lr}=2\times10^{-5}$ and batch size 16 on CPU using the \texttt{transformers} library \cite{wolf-etal-2020-transformers}.
We report macro F1, precision, and recall — standard for imbalanced hate speech evaluation.

\subsection{Results}

\begin{table}[t]
\centering
\small
\setlength{\tabcolsep}{4pt}
\begin{tabular}{lrrrr}
\toprule
\textbf{Condition} & \textbf{N} & \textbf{F1} & \textbf{Prec} & \textbf{Rec} \\
\midrule
\textsc{real-only}       & 500  & \textbf{0.715} & 0.715 & 0.715 \\
\textsc{real+top50}     & 1085 & \underline{0.727} & 0.741 & 0.730 \\
\midrule
\textsc{unfiltered}     & 1170 & 0.576 & 0.720 & 0.623 \\
\textsc{top-90\%}       & 1052 & 0.548 & 0.762 & 0.613 \\
\textsc{top-70\%}       &  818 & 0.575 & 0.697 & 0.618 \\
\textsc{top-50\%}       &  585 & 0.574 & 0.725 & 0.623 \\
\textsc{top-30\%}       &  350 & 0.576 & 0.644 & 0.605 \\
\textsc{top-10\%}       &  116 & 0.333 & 0.250 & 0.500 \\
\bottomrule
\end{tabular}
\caption{Results across all eight training conditions on the HatEval test set ($n_{\text{test}}=400$).
  \textbf{Bold}: best overall. \underline{Underline}: best augmented.}
\label{tab:results}
\end{table}

Table~\ref{tab:results} and Figure~\ref{fig:pareto} summarize all eight conditions.
Real-only training achieves F1\,=\,0.715 (baseline).
No synthetic-only condition approaches this score;
the best synthetic variant is \textsc{unfiltered} (F1\,=\,0.576),
leaving a gap of 0.139 F1 that quality filtering cannot close.

\begin{figure}[t]
  \centering
  \includegraphics[width=\columnwidth]{pareto_curve.pdf}
  \caption{Pareto curve: F1-macro as a function of the quality threshold applied to
    synthetic data. The red dashed line marks the real-data baseline (F1\,=\,0.715).
    The curve is effectively flat between top-90\% and top-30\%
    (range: 0.548--0.576), with no filtering threshold recovering the performance gap.}
  \label{fig:pareto}
\end{figure}

% ── 5. Analysis ──────────────────────────────────────────────
\section{Analysis}
\label{sec:analysis}

\subsection{The Flat Pareto Curve}
Figure~\ref{fig:pareto} shows no monotone relationship between quality threshold
and F1; the curve is effectively flat for thresholds 30\%--90\%
(F1 range: 0.548--0.576, spread of 0.028).
A one-way ANOVA across the five synthetic conditions (excluding the degenerate top-10\%)
yields $F = 511.05$, $p < 0.001$ (bootstrapped F1 distributions, $n_{\text{boot}}=500$).
While nominally significant, this result is a statistical artefact:
all five F1 point estimates are compressed into a band of 0.028 F1 units,
and the ANOVA significance is driven entirely by the bootstrap variance structure rather
than any substantive threshold effect.
Effect sizes confirm this: no pairwise comparison between thresholds
exceeds a meaningful difference in absolute F1.
Excluding the degenerate top-10\% condition (116 samples, F1 collapses to 0.333 due to
label imbalance in training), no threshold consistently outperforms any other.
The root cause is visible in Figure~\ref{fig:quality_dist}:
the discriminator quality score distributions for both hate and non-hate synthetic samples
concentrate heavily near zero, indicating the model confidently identifies all synthetic
text as synthetic.
With near-perfect AUC (1.00) on real vs.\ synthetic classification, the discriminator
lacks sufficient within-synthetic discriminative signal to rank quality usefully.

\begin{figure}[t]
  \centering
  \includegraphics[width=\columnwidth]{quality_score_dist.pdf}
  \caption{Quality score distributions for synthetic hate and non-hate samples.
    Both distributions concentrate near zero, indicating the discriminator classifies
    nearly all synthetic samples as clearly non-real.
    Filtering thresholds at 10\%, 30\%, 50\%, 70\%, and 90\% are shown;
    their overlap explains why threshold selection yields no meaningful performance gain.}
  \label{fig:quality_dist}
\end{figure}

\subsection{Precision-Recall Asymmetry}

\begin{table}[t]
\centering
\small
\setlength{\tabcolsep}{5pt}
\begin{tabular}{lrrr}
\toprule
\textbf{Condition} & \textbf{Prec} & \textbf{Rec} & \textbf{Gap (P$-$R)} \\
\midrule
\textsc{real-only}   & 0.715 & 0.715 & \phantom{$-$}0.000 \\
\textsc{real+top50} & 0.741 & 0.730 & \phantom{$-$}0.011 \\
\midrule
\textsc{unfiltered}  & 0.720 & 0.623 & \phantom{$-$}0.097 \\
\textsc{top-90\%}    & 0.762 & 0.613 & \phantom{$-$}0.149 \\
\textsc{top-70\%}    & 0.697 & 0.618 & \phantom{$-$}0.079 \\
\textsc{top-50\%}    & 0.725 & 0.623 & \phantom{$-$}0.103 \\
\textsc{top-30\%}    & 0.644 & 0.605 & \phantom{$-$}0.039 \\
\bottomrule
\end{tabular}
\caption{Precision-Recall gap per condition ($n_{\text{test}}=400$; top-10\% excluded as degenerate).
  Synthetic-only conditions show systematically larger P--R gaps (0.039--0.149)
  than real-data conditions (0.000--0.011), indicating a consistent generative bias.}
\label{tab:pr_gap}
\end{table}

Table~\ref{tab:pr_gap} reveals the key diagnostic finding:
all synthetic-only conditions (excluding the degenerate top-10\%)
exhibit substantially larger precision-recall gaps (0.039--0.149)
than real-data conditions (0.000--0.011).

This asymmetry exposes a systematic structural bias in LLM-generated hate speech.
The generator produces \emph{prototypical}, explicit forms of hate speech effectively,
resulting in high precision on clear-cut cases, but fails to replicate the distributional
breadth of borderline, implicit, and context-dependent instances present in human-annotated
data, causing a systematic recall deficit.
Pearson correlation between synthetic proportion and PR-gap is $r = 0.115$ ($p = 0.786$),
not statistically significant due to the compressed range of synthetic-proportion values,
but the qualitative pattern is unambiguous:
all four synthetic-only conditions with sufficient data (top-30\% through unfiltered)
exhibit gaps of 0.039--0.149 versus $< 0.012$ for both real-data conditions.
This pattern is consistent across all evaluated thresholds,
strongly indicating a \emph{generative bias} that is orthogonal to filtering threshold.

\begin{figure}[t]
  \centering
  \includegraphics[width=\columnwidth]{pr_asymmetry.pdf}
  \caption{Comparison of F1, precision, recall, and accuracy between REAL-ONLY and REAL+TOP50 augmented conditions (ntest = 400). Augmentation yields marginal gains of +1.7\% to +3.5\% across metrics, none reaching statistical significance.}
  \label{fig:pr_asymmetry}
\end{figure}

\subsection{Does Augmentation Help?}
The \textsc{real+top50} condition achieves F1\,=\,0.727, a nominally positive gain
of $+0.012$ over real-only (0.715).
McNemar's test on the 400-sample paired predictions yields
$\chi^2 = 3.125$, $p = 0.077$, falling short of the conventional $\alpha = 0.05$ threshold.
Bootstrapped Cohen's $d = 2.315$ [large], 95\%\,CI\,$[2.203, 2.419]$ (real vs.\ synthetic unfiltered)
confirms that the real-synthetic gap is large and replicable, while the real-augmented difference
remains marginal in practical terms. Note this effect size quantifies the real-versus-synthetic gap rather than the augmentation gain directly.
With 400 test samples, our design has adequate power to detect medium effects;
the non-significant McNemar result is therefore informative rather than underpowered.
We conclude that discriminator-filtered synthetic augmentation provides no statistically
or practically significant improvement, consistent with \cite{casula-etal-2024-delving}'s mixed findings.

% ── 6. Limitations ───────────────────────────────────────────
\section{Limitations}
\label{sec:limitations}

\paragraph{Single dataset and language.}
All experiments use English HatEval data from a single source distribution; results may be sensitive to annotation conventions and may not generalise across hate speech typologies or languages.

\paragraph{Quality proxy.}
The discriminator-based quality score is a proxy that measures distributional
closeness to real text, not semantic quality or label fidelity.
Human evaluation of synthetic outputs for annotation alignment, fluency, and
toxicity calibration remains out of scope and constitutes important future work.

\paragraph{Generator model and scale.}
LLaMA 3.1-8B via the Groq API may carry capability or stylistic biases
that shape the quality score distribution.
Larger models, instruction-tuned generators specifically targeting borderline cases,
or models fine-tuned on hate speech corpora may produce qualitatively different
synthetic data and alter the Pareto frontier.

\paragraph{Single quality metric.}
Our Pareto sweep uses one discriminator-based score.
Alternative quality signals—such as perplexity, BERTScore, or label-consistency filters—
may capture complementary aspects of sample quality and are worth exploring.

% ── 7. Conclusion ────────────────────────────────────────────
\section{Conclusion}
\label{sec:conclusion}

We presented the first systematic Pareto analysis of quality-filtered LLM-generated
hate speech data, sweeping six continuous quality thresholds with a fully open-source pipeline.
Evaluated on a 400-sample held-out test set, quality filtering does not close the gap
between synthetic and real data: the best synthetic-only condition achieves F1\,=\,0.576
against 0.715 for real-only training, a deficit of 0.139 F1 that no threshold recovers.
A precision-recall asymmetry analysis explains the failure mechanistically:
LLM-generated hate speech biases classifiers toward prototypical, explicit patterns
(high precision) while systematically suppressing recall of borderline and implicit cases.
Augmentation with filtered synthetic data yields F1\,=\,0.727 ($+0.012$), which
falls short of statistical significance (McNemar $\chi^2 = 3.125$, $p = 0.077$).

Our practical recommendation is clear: LLM-generated hate speech data should be treated
as complementary augmentation with calibrated expectations, not a substitute
for human annotation, particularly in domains where borderline cases drive evaluation.
Future work should investigate targeted generators for implicit hate speech,
human-in-the-loop quality evaluation, and cross-lingual transfer of the filtering pipeline.

% ── References ───────────────────────────────────────────────
\bibliography{references}
\bibliographystyle{plain}

\end{document}